\VerdierPhysicalPageAnchor{0062}%
\VerdierSemanticAnchor{ch1-1-6-4}
\VerdierSemanticAnchor{ch1-1-6-4-definition}
\noindent\textbf{Définition 1.6.4.}\quad
On appelle \emph{cocycle de Koszul relatif à une base ordonnée
$e_1,\ldots,e_n$ de $\mathbb Z^n$} l'unique cocycle $\alpha$ dans la
classe de Koszul relative à la base $e_1,\ldots,e_n$ qui vérifie les
égalités :

\VerdierSemanticAnchor{ch1-1-6-4-eq1}
\begin{equation}
  \alpha(\sigma_1,\sigma_2)=1
  \quad
  \forall\sigma_1\in\mathbb Z^n(\leq i),\quad
  \forall\sigma_2\in\mathbb Z^n(\geq i),\quad
  \forall i\in[1,n].
  \tag{1.6.4.1}
\end{equation}

\VerdierSemanticAnchor{ch1-1-6-4-coordinate-intro}
Soient
$\sigma_1=\sum_{i=1}^{n}p_i e_i$,
$\sigma_2=\sum_{i=1}^{n}q_i e_i$
deux éléments de $\mathbb Z^n$ ; on a alors :

\VerdierSemanticAnchor{ch1-1-6-4-eq2}
\begin{equation}
  \alpha(\sigma_1,\sigma_2)
  =(-1)^{\left(\sum_{i<j}q_i p_j\right)}.
  \tag{1.6.4.2}
\end{equation}

\VerdierSemanticAnchor{ch1-1-6-4-justification}
Cette définition demande une justification. L'existence d'un cocycle
dans la classe de Koszul vérifiant (1.6.4.1) est claire : un calcul
immédiat montre que le cocycle (1.6.4.2) vérifie (1.6.4.1) et qu'il est
dans la classe de Koszul (1.6.2). Pour vérifier l'unicité, il suffit de
remarquer que tout cocycle cohomologue au cocycle $\alpha$ est du type :

\VerdierSemanticAnchor{ch1-1-6-4-cohomologous-form}
\begin{equation*}
  \alpha_1(\sigma_1,\sigma_2)
  =\alpha(\sigma_1,\sigma_2)
   \beta(\sigma_2)\beta(\sigma_1+\sigma_2)^{-1}\beta(\sigma_1)
\end{equation*}
et qu'un tel cocycle possède la propriété (1.6.4.1) si et seulement si :

\VerdierSemanticAnchor{ch1-1-6-4-condition}
\begin{equation*}
  \beta(\sigma_1+\sigma_2)=\beta(\sigma_1)\beta(\sigma_2)
  \quad
  \forall i,\quad
  \forall\sigma_1\in\mathbb Z^n(\leq i),\quad
  \forall\sigma_2\in\mathbb Z^n(\geq i).
\end{equation*}

\VerdierSemanticAnchor{ch1-1-6-4-uniqueness}
On en déduit immédiatement que
$\beta\colon\mathbb Z^n\longrightarrow\mathbb Z/2\mathbb Z$
est un homomorphisme de groupes et par suite qu'on a
$\alpha_1=\alpha$.

\VerdierSemanticAnchor{ch1-1-6-5}
\VerdierSemanticAnchor{ch1-1-6-5-order-dependence}
\noindent\textbf{1.6.5.}\quad
La définition du cocycle de Koszul fait intervenir de façon essentielle
l'ordre total que l'on met sur la base $e_1,\ldots,e_n$ de
$\mathbb Z^n$. Désignons par $\leq$ l'ordre naturel de $[1,n]$ et par
$\leq_s$ un ordre total sur $[1,n]$. Soit $s$ l'unique permutation de
$[1,n]$ telle que :

\VerdierSemanticAnchor{ch1-1-6-5-eq1}
\begin{equation}
  i\leq_s j\quad\Longleftrightarrow\quad s(i)\leq s(j).
  \tag{1.6.5.1}
\end{equation}

\VerdierSemanticAnchor{ch1-1-6-5-cocycles}
Soient $\alpha$ et $\alpha_s$ les cocycles de Koszul relatifs
respectivement aux ordres $\leq$ et $\leq_s$. On a alors :

\VerdierSemanticAnchor{ch1-1-6-5-eq2}
\VerdierSemanticAnchor{ch1-1-6-5-eq2-natural}
\VerdierSemanticAnchor{ch1-1-6-5-eq2-coordinates}
\VerdierSemanticAnchor{ch1-1-6-5-eq2-permuted}
\begin{equation}
  \begin{aligned}
    \alpha(\sigma_1,\sigma_2)
      &=(-1)^{\left(\sum_{i<j}q_i p_j\right)},\\
    \sigma_1&=\sum_i p_i e_i,
      \qquad \sigma_2=\sum_i q_i e_i,\\
    \alpha_s(\sigma_1,\sigma_2)
      &=(-1)^{\left(\sum_{s(i)<s(j)}q_i p_j\right)},
  \end{aligned}
  \tag{1.6.5.2}
\end{equation}
\VerdierSemanticAnchor{ch1-1-6-5-open}
